

\documentclass[11pt]{article}
\usepackage{amsmath}
\usepackage{hyperref}
\usepackage[
backend=biber,
style=numeric,
sorting=ynt
]{biblatex}
\addbibresource{report2bib.bib}

\title{C\# Implementation of the Fast Fourier Transform (FFT) Algorithm}
\author{Ike Halderan, Josiah Hoppe, Jacari Jellum, Sauharda Shrestha}
\date{\today}

\begin{document}
\maketitle

%\begin{abstract}
%This is a short abstract summarizing your article.
%\end{abstract}

\section*{Introduction}

\section*{Mathematical Validity}
We will first show that any function approximated at $2n$ points can be approximated by a Fourier series, by reproducing the approach shown in \cite{hamming}, while using the same definitions as \cite{textbook}. The set of functions 
\[1, \cos x, \cos 2x, \ldots, \cos \left(n-1\right)x, \cos nx\]
\[\sin x, \sin 2x, \dots, \sin \left(n-1\right)x, \sin nx\]

is orthogonal over the discrete set of points 

\[0, \frac{\pi}{n}, \frac{2\pi}{n}, \ldots, \frac{\left(2n-1\right)\pi}{n}\]

The orthogonality relations, when $k \leq n, m \leq n$, are

\begin{align*}
    \sum_{x=0}^{2n-1} \sin \frac{\pi}{n}kx \sin \frac{\pi}{n}mx & = \begin{cases}
        0 & \text{if } k \neq m \\
        n & \text{if } k = m \neq 0 \\        
    \end{cases} \\
    \sum_{x=0}^{2n-1} \sin \frac{\pi}{n}kx \cos \frac{\pi}{n}mx & = 0 \\
    \sum_{x=0}^{2n-1} \cos \frac{\pi}{n}kx \cos \frac{\pi}{n}mx & = \begin{cases}
        0 & \text{if } k \neq m \\
        n & \text{if } k = m \neq 0 \\
    \end{cases}
\end{align*}

Suppose $f(x)$ is an arbitrary function. Since each of the above functions is orthogonal to each other, we can write $f(x)$ as a linear combination of these functions, so

\[f(x) = \frac{a_0}{2} + \sum_{k=1}^{n-1} \left( a_k \cos \frac{\pi}{n}kx + b_k \sin \frac{\pi}{n}kx \right) + \frac{a_n}{2} \cos \frac{\pi}{n}nx.\]

Then, we can use the orthogonality relations to solve for the coefficients $a_k$ and $b_k$. For example, to solve for $a_k$, we multiply both sides of the equation by $\cos \frac{\pi}{n}mx$ and sum over all $x$ from $0$ to $2n-1$. Thus, we have


\section*{Algorithm Overview}


\section*{Conclusion}


\printbibliography[
heading=bibintoc,
title={Whole bibliography}
]

\end{document}
